\documentclass[11pt,reqno]{amsart}
\usepackage[T1]{fontenc}
\usepackage{iftex}
\ifPDFTeX
  \usepackage[utf8]{inputenc}
\fi
\usepackage{lmodern}
\usepackage{amsmath,amssymb,amsthm,mathtools,mathrsfs}
\usepackage{booktabs}
\usepackage{enumitem}
\usepackage{microtype}
\usepackage{xcolor}
\usepackage[hidelinks]{hyperref}
\usepackage[nameinlink,noabbrev]{cleveref}

\numberwithin{equation}{section}
\raggedbottom

\newtheorem{theorem}{Theorem}[section]
\newtheorem{proposition}[theorem]{Proposition}
\newtheorem{lemma}[theorem]{Lemma}
\newtheorem{corollary}[theorem]{Corollary}
\newtheorem{question}[theorem]{Question}
\theoremstyle{definition}
\newtheorem{definition}[theorem]{Definition}
\newtheorem{example}[theorem]{Example}
\theoremstyle{remark}
\newtheorem{remark}[theorem]{Remark}

\newcommand{\A}{\mathbb A}
\newcommand{\C}{\mathbb C}
\newcommand{\F}{\mathbb F}
\newcommand{\Gm}{\mathbb G_{\mathrm m}}
\newcommand{\Lef}{\mathbb L}
\newcommand{\PP}{\mathbb P}
\newcommand{\Cl}{\operatorname{Cl}}
\newcommand{\Disc}{\operatorname{Disc}}
\newcommand{\Res}{\operatorname{Res}}
\newcommand{\Sing}{\operatorname{Sing}}
\newcommand{\Sym}{\operatorname{Sym}}
\newcommand{\Spec}{\operatorname{Spec}}
\newcommand{\ord}{\operatorname{ord}}
\newcommand{\id}{\operatorname{id}}
\newcommand{\cE}{E_{\mathrm c}}
\newcommand{\cchi}{\chi_{\mathrm c}}
\newcommand{\set}[1]{\left\{#1\right\}}
\newcommand{\abs}[1]{\left|#1\right|}
\newcommand{\angles}[1]{\left\langle#1\right\rangle}

\crefname{theorem}{theorem}{theorems}
\Crefname{theorem}{Theorem}{Theorems}
\crefname{proposition}{proposition}{propositions}
\Crefname{proposition}{Proposition}{Propositions}
\crefname{lemma}{lemma}{lemmas}
\Crefname{lemma}{Lemma}{Lemmas}
\crefname{corollary}{corollary}{corollaries}
\Crefname{corollary}{Corollary}{Corollaries}
\crefname{question}{question}{questions}
\Crefname{question}{Question}{Questions}


\title[Filtered rigidity]{Filtered Rigidity of the Degree-Seven Jacobian
Counterexample:\\A Transverse Algebra of Length \(584\)}
\author{Nathaniel Monson}
\date{July 22, 2026}

\begin{document}

\begin{abstract}
Let \(G\colon\A^3_{\C}\to\A^3_{\C}\) be the normalized degree-seven
counterexample to the Jacobian conjecture, and let \(K_{3,7}\) be the
finite-dimensional coefficient scheme of maps \(H\) satisfying
\[
\deg H\le7,\qquad H(0)=0,\qquad JH(0)=I,\qquad \det JH=1.
\]
Unrestricted Keller deformations of \(G\) are volume-preserving formal source
reparametrizations.  Ordinary degree is not preserved by those
reparametrizations, and the bounded-degree scheme cuts the formal coordinate
orbit nontransversely.  After removing the normalized affine source action,
an explicit torus-stable affine slice has a ten-dimensional tangent
space, with weights
\[
(-1,2,-3,-2,-1,0,1,1,2,3).
\]
Nevertheless, its reduced completed local ring is \(\C\).

We determine the entire nonreduced thickening by exact computer-assisted
calculation.  Its length is \(584\), its Hilbert--Samuel function is
\[
(1,10,44,108,157,145,86,30,3),
\]
and its maximal ideal satisfies \(\mathfrak m^9=0\ne\mathfrak m^8\).
Equivalently, the normalized affine orbit is the reduced local component and
occurs with transverse multiplicity \(584\).  The Macaulay inverse system has
60 minimal contraction generators, while the
Kuranishi ideal has 36 minimal equations: 11 quadratic, 13 cubic, 11
quartic, and one sextic.  In particular, the local algebra has
Cohen--Macaulay type 60 and is neither a complete intersection, Gorenstein,
nor level.  We state the exact specialization and rational-reconstruction
certificates used in the proof and distinguish this bounded-degree,
affine-quotiented theorem from unrestricted deformation claims.
\end{abstract}

\maketitle

\begin{center}
\small\emph{Working computer-assisted manuscript.  The exact calculations
have been rerun from the recovered source artifacts.  A genuinely independent
CAS implementation and specialist review remain submission requirements.}
\end{center}

\section{Introduction}
\label{sec:introduction}

Put \(u=1+xy\), and define
\begin{align}
P&=u^3z+y^2u(4+3xy),\label{eq:P}\\
Q&=y+3xu^2z+3xy^2(4+3xy),\label{eq:Q}\\
R&=2x-3x^2y-x^3z.\label{eq:R}
\end{align}
The announced counterexample is \(F=(P,Q,R)\); it has
\(\det JF=-2\) and is not injective
\cite{alpoge2026announcement,tao2026digestion,ulam2026counterexample}.
We use the normalized map
\begin{equation}
\label{eq:G}
G=(R/2,Q,P).
\end{equation}

\emph{Discovery attribution.}  We credit Akhil Mathew as the primary human
source of the problem: Mathew prompted Levent Alp\"oge, Alp\"oge prompted
Fable, Fable produced the example, and Alp\"oge announced it publicly
\cite{ulam2026counterexample}.
Then
\[
G(0)=0,\qquad JG(0)=I,\qquad \det JG=1,\qquad \deg G=7.
\]

The disproof immediately raises a deformation question: is this map an
isolated accident, or is it one point of a family?  The question has no useful
answer until both the allowed coordinate changes and a degree bound are
specified.  There are already unrestricted families of nonproper Keller maps
\cite{ulam2026counterexample,gallagher2026explained}.  The theorem in this
paper is instead local, bounded-degree, and transverse to a stated affine
action.

Let \(K_{3,7}\) be the coefficient scheme of polynomial maps
\(H\colon\A^3\to\A^3\) of ordinary degree at most seven satisfying
\begin{equation}
\label{eq:normalization}
H(0)=0,\qquad JH(0)=I,\qquad \det JH=1.
\end{equation}
The normalized affine source action is
\begin{equation}
\label{eq:rho}
\rho_{a,A}(H)(x)
  =[JH(a)A]^{-1}\bigl(H(a+Ax)-H(a)\bigr),
\qquad (a,A)\in\operatorname{Aff}_3.
\end{equation}
We construct an explicit affine slice \(S\subset K_{3,7}\) through \(G\).
Our main result describes
\[
\mathcal R=\widehat{\mathcal O}_{S,G}.
\]
The ring and all of its invariants may be formed over \(\mathbb Q\); the
statements below are unchanged after base extension to \(\C\).

\begin{theorem}[Transverse local algebra]
\label{thm:main}
For the slice \(S\) of \cref{def:slice}, the following hold.
\begin{enumerate}[label=(\roman*)]
\item The reduced completed local ring is a point:
\[
\mathcal R_{\mathrm{red}}\cong\C.
\]
\item If \(\mathfrak m\) is the maximal ideal, then
\[
\operatorname{length}\mathcal R=584,\qquad
\mathfrak m^9=0\ne\mathfrak m^8,
\]
and
\[
\bigl(\dim_\C\mathfrak m^d/\mathfrak m^{d+1}\bigr)_{d=0}^8
=(1,10,44,108,157,145,86,30,3).
\]
\item The Macaulay inverse system of \(\mathcal R\) has 60 minimal
contraction generators.  Their top divided-power degrees are distributed as
\[
\begin{array}{c|rrrr}
\text{degree}&5&6&7&8\\ \hline
\text{number}&2&33&22&3.
\end{array}
\]
Consequently
\(\dim_\C\operatorname{Soc}(\mathcal R)=60\).
\item The Kuranishi ideal has 36 minimal generators, distributed by initial
order as
\[
\begin{array}{c|rrrr}
\text{order}&2&3&4&6\\ \hline
\text{number}&11&13&11&1.
\end{array}
\]
For the minimal resolution over
\(\C[[u_1,\ldots,u_{10}]]\), the extremal Betti numbers are
\(\beta_1=36\) and \(\beta_{10}=60\).
\end{enumerate}
\end{theorem}

The tangent space is not zero: it has dimension ten.  Thus the content of
\cref{thm:main} is not infinitesimal rigidity.  Every first-order transverse
direction is obstructed, and all of the resulting infinitesimal thickness is
captured by a comparatively large Artin algebra.

\subsection*{Relation to concurrent work}

The mixed-sign grading of the map and its geometric significance were
already public \cite{shaska2026graded}.  Teller computed a residual tangent
dimension of ten and modular obstructions through finite order
\cite{teller2026rigidity}.  Mayner's repository contains an exact deformation
calculation in a narrower fixed-direction, fixed-support ansatz; its
provenance statement credits the mathematical work to Claude Fable 5 and the
circulation and packaging to William G.\ P.\ Mayner
\cite{mayner2026structural}.  These are important predecessors, but neither
determines the reduced characteristic-zero germ of the full normalized
degree-seven slice or its complete local algebra.

Our candidate contribution begins with the exact residual weight
representation and continues through the characteristic-zero radical
statement, length, inverse system, socle, and minimal equations.  We make no
claim to the tangent dimension alone, qualitative nonreducedness, the torus
action itself, or rigidity in a restricted ansatz.

\subsection*{Formal orbit versus bounded degree}

A Keller map is \'{e}tale.  Formal \'{e}taleness gives unique lifts, so every
unrestricted normalized Keller deformation comes from a volume-preserving
formal change of source coordinates \cite{stacks02HF}.  This standard fact
does not imply the present theorem.  The formal changes need not preserve
ordinary degree at most seven.  The geometry below is the high-order
incidence of that formal orbit with the bounded-degree coefficient space.

\begin{proposition}[Formal source-orbit theorem]
\label{prop:formal-orbit}
Let \(A\) be a local Artin \(\C\)-algebra, and let \(H_A\) be a normalized
Keller deformation of \(G\) over \(A\) reducing to \(G\).  There is a unique
normalized formal \(A\)-automorphism of \(\widehat{\A}^3_0\),
\(\phi_A\equiv\id\pmod{\mathfrak m_A}\) such that
\[
H_A=G\circ\phi_A.
\]
Moreover, \(\det J\phi_A=1\).  Consequently the degree-\(\le7\)
deformation functor is the subfunctor of volume-preserving source
reparametrizations cut out by
\[
\deg(G\circ\phi_A)\le7.
\]
\end{proposition}

\begin{proof}
On the completed local source, formal \'{e}taleness makes \(G_A\) a formal
isomorphism.  Set \(\phi_A=G_A^{-1}\circ H_A\).  It is the unique formal
automorphism reducing to the identity and satisfying
\(G_A\circ\phi_A=H_A\).  Finally,
\[
1=\det JH_A=\det JG(\phi_A)\det J\phi_A=\det J\phi_A.
\]
The normalizations at the origin follow from uniqueness.
\end{proof}

\section{The affine action and a transverse slice}
\label{sec:slice}

The coefficient affine space underlying \eqref{eq:normalization}, before the
determinant equations are imposed, has dimension
\[
3\left(\binom{10}{3}-4\right)=348.
\]
The map \(G\) is equivariant for the source torus
\begin{equation}
\label{eq:torus}
A_\tau=\operatorname{diag}(\tau^{-1},\tau,\tau^2),
\qquad
G(A_\tau x)=A_\tau G(x).
\end{equation}

\begin{proposition}[Normalized affine stabilizer]
\label{prop:stabilizer}
The stabilizer of \(G\) under \eqref{eq:rho} is
\[
\operatorname{Stab}_{\operatorname{Aff}_3}(G)
=\set{A_\tau:\tau\in\Gm}\cong\Gm.
\]
\end{proposition}

\begin{proof}
Suppose \(\rho_{a,A}(G)=G\).  With
\[
\alpha(x)=a+Ax,\qquad
\beta(v)=G(a)+JG(a)Av,
\]
one has \(G\circ\alpha=\beta\circ G\).  The omitted locus of \(G\) is
\[
\Gamma=\set{\left(\frac{s}{3},\frac2s,\frac1{3s^2}\right):s\in\C^\times}
=V(3v_1v_2-2,\,27v_1^2v_3-1)
\cong\Gm
\]
\cite{ulam2026counterexample}.  Hence \(\beta\) preserves \(\Gamma\).
Restrictions of affine-linear functions to this parametrization span
\[
\angles{1,s,s^{-1},s^{-2}}\subset\C[s,s^{-1}].
\]
The inversion \(s\mapsto\mu/s\) does not preserve this space, whereas
\(s\mapsto\mu s\) does.  No nonzero affine-linear function vanishes on
\(\Gamma\), so the ambient affine extension is unique and equals
\[
\operatorname{diag}(\mu,\mu^{-1},\mu^{-2}).
\]

It remains to identify the source lift.  The generic fiber is described by
\[
cv^3-2v^2+bv-2a=0.
\]
This cubic is irreducible and has nonsquare discriminant
\[
-4(27a^2c^2-18abc+16a+b^3c-b^2);
\]
For irreducibility, put \(K=\C(b,c)\) and use the valuation
\(\nu_\infty\) of \(K(a)\) at \(a=\infty\), normalized by
\(\nu_\infty(a)=-1\).  If a root \(r\in K(a)\) existed and
\(m=\nu_\infty(r)\), then for \(m\ge0\) the term \(-2a\) would have the
unique least valuation.  For \(m<0\), the four term valuations would be
\[
3m,\qquad 2m,\qquad m,\qquad -1.
\]
Cancellation at the least valuation would require \(3m=-1\), impossible
for integral \(m\).  Thus the cubic has no root in \(K(a)\), hence is
irreducible.  The discriminant is nonsquare because the discriminant of the
parenthesized expression as a quadratic in \(a\) is
\(-4(3bc-4)^3\).  The normal closure therefore has group \(S_3\).
For the corresponding subgroup \(S_2\subset S_3\),
\[
N_{S_3}(S_2)/S_2=1,
\]
so the cubic function-field extension has no nontrivial automorphism over
the target field.  The known source lift is \(A_{\mu^{-1}}\), and uniqueness
forces \(\alpha=A_{\mu^{-1}}\).  Finally, finite-type group schemes in
characteristic zero are smooth, so no nonreduced stabilizer remains.
\end{proof}

\begin{definition}[The slice]
\label{def:slice}
For a perturbation \(H-G\), set the coefficients of the following eleven
component--monomial pairs to zero:
\[
\begin{array}{c|l}
\text{component}&\text{monomials}\\ \hline
1&y^2,\ xz,\ xy,\ x^2,\ x^3\\
2&y^2,\ xz,\ xy,\ x^2\\
3&y^2,\ xy.
\end{array}
\]
Let \(A_{\mathrm{sl}}\) be the resulting 337-dimensional affine subspace and
put
\[
S=K_{3,7}\cap A_{\mathrm{sl}}.
\]
\end{definition}

\begin{proposition}[Transversality]
\label{prop:transversality}
The slice is torus-stable.  The selected \(11\times11\) affine-orbit minor
has determinant \(-23328\).  In particular,
\[
\operatorname{Aff}_3\times^{\Gm}S\longrightarrow K_{3,7}
\]
is \'{e}tale at \([1,G]\).
\end{proposition}

\begin{proof}
The eleven coefficient perturbations are weight vectors for
\eqref{eq:torus}; the two whose base coefficients are nonzero have weight
zero.  This proves stability.

First work in the smooth \(348\)-dimensional ambient coefficient space
\(\mathcal A\).  The action map
\[
\operatorname{Aff}_3\times^{\Gm}A_{\mathrm{sl}}\longrightarrow\mathcal A
\]
has source and target of the same dimension.  Differentiating
\eqref{eq:rho} in the twelve affine parameters, quotienting the
one-dimensional stabilizer, and applying the eleven displayed coefficient
functionals gives determinant \(-23328\).  The ambient action map is
therefore \'{e}tale at \([1,G]\).

The Keller scheme \(K_{3,7}\subset\mathcal A\) is invariant under the
normalized affine action.  Hence the inverse image of \(K_{3,7}\) under the
ambient action map is, near \([1,G]\),
\[
\operatorname{Aff}_3\times^{\Gm}(K_{3,7}\cap A_{\mathrm{sl}}).
\]
The claimed morphism is thus the base change of the ambient \'{e}tale map
along \(K_{3,7}\hookrightarrow\mathcal A\), and is \'{e}tale.
\end{proof}

\begin{corollary}[Local product and multiplicity]
\label{cor:local-product}
There is a noncanonical formal isomorphism
\[
\widehat{\mathcal O}_{K_{3,7},G}
\cong
\C[[t_1,\ldots,t_{11}]]
\widehat\otimes_\C\mathcal R.
\]
Once \(\mathcal R_{\mathrm{red}}=\C\) is proved below, the reduced local
component is the normalized affine orbit and its scheme-theoretic
transverse multiplicity is
\[
\operatorname{length}\mathcal R=584.
\]
The formal local quotient stack is
\[
[\operatorname{Spf}\mathcal R/\Gm].
\]
\end{corollary}

\begin{proof}
The \'{e}tale morphism of \cref{prop:transversality} identifies the completed
local germ with the product of the eleven-dimensional orbit direction and
the slice, with the residual stabilizer \(\Gm\).  Length of the transverse
Artin factor is the component multiplicity.
\end{proof}

\section{The equivariant Kuranishi map}
\label{sec:kuranishi}

Let
\[
\Phi(H)=\det JH-1.
\]
The differential \(D\Phi_G\) has rank \(327\) on the 337-dimensional slice.
Choose tangent coordinates \(u_1,\ldots,u_{10}\) by taking the following free
coefficient functionals; component indices in the table are one-based.
\[
\begin{array}{c|c|r}
&\text{free coefficient}&\text{weight}\\ \hline
u_1&H_3[x^2y^3]&-1\\
u_2&H_3[xy^5]&2\\
u_3&H_3[x^4y^3]&-3\\
u_4&H_3[x^4y^2z]&-2\\
u_5&H_3[x^3y^4]&-1\\
u_6&H_3[x^3y^3z]&0\\
u_7&H_3[x^3y^2z^2]&1\\
u_8&H_3[x^2y^5]&1\\
u_9&H_3[x^2y^4z]&2\\
u_{10}&H_3[xy^6]&3.
\end{array}
\]
Formal implicit elimination of the 327 pivot coefficient directions gives a
torus-equivariant Kuranishi map
\[
\kappa\colon\widehat{\A}^{10}\longrightarrow
\operatorname{coker}(D\Phi_G)
\]
and an isomorphism
\begin{equation}
\label{eq:kuranishi-ring}
\mathcal R\cong
\C[[u_1,\ldots,u_{10}]]/I_\kappa.
\end{equation}
The exact quadratic part spans an eleven-dimensional space.  In particular,
one equation begins with \((u_1-3u_5)^2\), but the quadratic cone is
positive-dimensional; the reduced-rigidity theorem is not a consequence of
the quadratic approximation.

\section{Reduced local rigidity}
\label{sec:reduced}

We first isolate the geometric lemma that converts three smaller exact
calculations into a statement about the full germ.

\begin{lemma}[Torus nullcone]
\label{lem:nullcone}
Let \(X\) be an affine finite-type \(\Gm\)-scheme with a fixed point \(0\),
equivariantly embedded in
\[
V=V_-\oplus V_0\oplus V_+.
\]
If the reduced germs at zero of
\[
X\cap V_-,\qquad X\cap V_0,\qquad X\cap V_+
\]
are all zero-dimensional, then the reduced germ
\((X_{\mathrm{red}},0)\) is zero-dimensional.
\end{lemma}

\begin{proof}
Suppose a positive-dimensional irreducible component
\(Y\subset X_{\mathrm{red}}\) passes through zero.  The connected group
\(\Gm\) preserves \(Y\).  If its action on \(Y\) is trivial, then
\(Y\subset V_0\), a contradiction.

Otherwise the generic orbit has dimension one, and
\[
\dim(Y/\!/\Gm)\le\dim Y-1.
\]
The fiber of \(Y\to Y/\!/\Gm\) over the image of zero therefore has positive
local dimension at zero.  Let \(Z\) be a positive-dimensional irreducible
component of that fiber through zero.

Every invariant regular function has on \(Z\) its value at zero.  Thus all
weight-zero coordinate functions vanish on \(Z\).  If a positive-weight
coordinate \(x_i\) and a negative-weight coordinate \(y_j\) were both
nonzero at a point of \(Z\), the balanced monomial
\[
x_i^{|\operatorname{wt}(y_j)|}
y_j^{\operatorname{wt}(x_i)}
\]
would be a nonzero invariant there, a contradiction.  Hence
\(Z\subset V_+\cup V_-\).  Irreducibility puts \(Z\) in one of these two
subspaces, contradicting the hypotheses.
\end{proof}

\begin{proposition}[Exact pure-weight certificates]
\label{prop:pure-weight}
For the Kuranishi germ of \eqref{eq:kuranishi-ring}, the following hold.
\begin{enumerate}[label=(\roman*)]
\item On the positive locus, with parameters
\[
(p_1,p_2,p_3,p_4,p_5)=(u_7,u_8,u_2,u_9,u_{10}),
\]
the eliminated ideal contains
\[
p_1^6,\quad p_2^6,\quad p_3^3,\quad p_4^3,\quad p_5^2.
\]
\item On the negative locus, with parameters
\[
(n_1,n_2,n_3,n_4)=(u_1,u_5,u_4,u_3),
\]
the eliminated ideal contains
\[
n_1^4,\quad n_2^4,\quad n_3^4,\quad n_4^3.
\]
\item The weight-zero locus has one tangent parameter \(q=u_6\).  Its first
nonzero compatibility equation occurs in order three and has coefficient
\[
\frac{212135552}{304438725}q^3.
\]
\end{enumerate}
\end{proposition}

\begin{proof}[Computational proof]
The determinant equations are triangular by torus weight.  On the positive
locus, variables in the current weight block enter through \(D\Phi_G\),
while nonlinear terms use lower positive weights.  Exact recursive
elimination through weight six leaves five free tangent parameters; a
rational Gr\"obner basis has 24 elements and contains the five displayed
powers.  The negative calculation reverses the weights.  Elimination through
absolute weight eight produces a 16-element rational Gr\"obner basis
containing the four displayed powers.

In weight zero there are 21 coefficient variables and the linearized
determinant map has rank 20.  Solving the pivot variables formally in \(q\)
produces no order-two compatibility condition.  At order three, the
coefficient of \(x^8y^4z^2\) is the displayed nonzero multiple of \(q^3\).
The scripts, exact outputs, and hash-pinned reproduction record are described
in \cref{sec:computation}.
\end{proof}

\begin{theorem}[Reduced affine rigidity]
\label{thm:reduced}
The reduced completed slice is a point:
\[
\mathcal R_{\mathrm{red}}\cong\C.
\]
Equivalently,
\[
\sqrt{I_\kappa}=(u_1,\ldots,u_{10}).
\]
\end{theorem}

\begin{proof}
The slice \(S\) is finite type.  The formal implicit elimination
\eqref{eq:kuranishi-ring} identifies its completed local ring at \(G\) with
the displayed Kuranishi quotient.  For each pure-weight coordinate
subspace \(V_-\), \(V_0\), and \(V_+\), completion commutes with the
corresponding closed intersection.  By \cref{prop:pure-weight}, each such
completed local intersection is Artinian.  Faithful flatness of completion,
and preservation of Krull dimension for Noetherian local rings, therefore
make each finite-type local germ \((S\cap V_\epsilon,G)\) zero-dimensional.

Apply \cref{lem:nullcone} to the finite-type torus-stable slice.  The full
reduced germ is zero-dimensional.  Since its support contains the base point
and the radical is a homogeneous maximal ideal in the completed tangent
coordinates, the displayed equality follows.
\end{proof}

\begin{remark}
\label{rem:scope}
The theorem says that there is no nonconstant reduced branch through \(G\)
inside \(K_{3,7}\), transverse to the normalized affine action.  It does not
say that \(G\) has no degree-increasing deformation, no deformation before
quotienting, or no deformation in a stabilized dimension.
\end{remark}

\section{Length, Hilbert function, and Loewy length}
\label{sec:length}

Write \(S_0=\mathbb Q[[u_1,\ldots,u_{10}]]\), let
\(R=S_0/I_\kappa\), and let \(\mathfrak m\) denote its maximal ideal.  The
characteristic-zero computation is made over \(\mathbb Q\); base change gives
the complex ring of the preceding sections.

For \(N\ge0\), put
\[
D_{\le N}=
\left(R/\mathfrak m^{N+1}\right)^\vee.
\]
We use the divided-power convention, so multiplication by \(u_i\) is dual to
contraction
\[
\partial_iY^{[\alpha]}=Y^{[\alpha-e_i]}.
\]

\begin{proposition}[The three terminal classes]
\label{prop:terminal}
The order-seven quotient has dimension \(581\).  The degree-eight dual layer
has dimension three, with one class in each torus weight \(-1,0,1\).
Consequently
\[
\dim_\mathbb Q R/\mathfrak m^9=584
\]
and
\[
\operatorname{ch}(\mathfrak m^8/\mathfrak m^9)=z^{-1}+1+z.
\]
\end{proposition}

\begin{proof}[Computer-assisted proof]
Take the good prime \(p=1{,}000{,}003\).  The modular calculation gives
\[
\dim R/\mathfrak m^9\le584,
\]
with at most one new degree-eight class in each of the three displayed
weights.  Rank can only fall under this specialization.  The modular quotient
dimension is therefore an upper bound in characteristic zero.

For the reverse inequality, the exact Kuranishi recursion was evaluated on
the coordinatewise down-set of the 60 degree-eight tangent monomials that can
occur in the candidate classes.  This produced 7,300 exact coefficients
through order eight, including 990 in order eight.  Every coefficient through
order seven agrees with the full rational calculation, and all 990
order-eight coefficients agree modulo \(p\) with the independently generated
full modular output.

Chinese remaindering and rational reconstruction then give three rational
dual classes.  They have supports of sizes 192, 387, and 307 in weights
\(-1,0,1\), and satisfy respectively 3,045, 6,236, and 5,363 exact
interacting rows.  Each is normalized by a degree-eight coefficient equal to
one, and their images are independent modulo \(D_{\le7}\).  Thus the lower
bound is \(581+3=584\), proving equality.
\end{proof}

\begin{proposition}[Order-nine integrability certificate]
\label{prop:order-nine}
One has
\[
\mathfrak m^9=0.
\]
\end{proposition}

\begin{proof}
Let \(\Lambda^{(9)}\) be a hypothetical nonzero top-degree-nine dual class of
torus weight \(W\).  Every contraction \(\partial_i\Lambda^{(9)}\) lies in the
three-dimensional top degree-eight dual space.  Hence
\[
\partial_i\Lambda^{(9)}
=a_i\Lambda_{W-w_i}^{(8)}
\]
when \(W-w_i\in\{-1,0,1\}\), and it is zero otherwise.  Equality of mixed
contractions produces a homogeneous linear system in the active scalars
\(a_i\).

For each possible weight \(W=-4,\ldots,4\), an exact square subsystem has the
following determinant.
\[
\begin{array}{c|r|r}
W&\text{active scalars}&\text{determinant}\\ \hline
-4&1&248832\\
-3&2&23219011584\\
-2&4&159739999685311463424\\
-1&4&59902499881991798784\\
0&5&4416491511299491340746752\\
1&5&-14905658850635783275020288\\
2&5&-1656184316737309251780032\\
3&3&-71328803586048\\
4&1&-27648
\end{array}
\]
All are nonzero.  For \(|W|>4\), there are no active contractions.  Thus every
\(a_i\) is zero and all contractions of \(\Lambda^{(9)}\) vanish, which forces
\(\Lambda^{(9)}=0\).  Therefore
\(\mathfrak m^9/\mathfrak m^{10}=0\).  Nakayama's lemma applied to the
finitely generated module \(\mathfrak m^9\) gives \(\mathfrak m^9=0\).
\end{proof}

\begin{corollary}[Hilbert--Samuel function]
\label{cor:hilbert}
The ring \(R\) has length \(584\), Loewy length nine, and Hilbert--Samuel
function
\[
(1,10,44,108,157,145,86,30,3).
\]
\end{corollary}

\begin{proof}
The exact ranks of the successive filtered quotients through degree seven are
\[
(1,10,44,108,157,145,86,30).
\]
\Cref{prop:terminal} adds the terminal dimension three, and
\cref{prop:order-nine} proves there is no further layer.  The sum is \(584\).
\end{proof}

\section{The complete inverse system}
\label{sec:inverse}

Let
\[
D=I_\kappa^\perp
\subset\mathbb Q_{\mathrm{DP}}[Y_1,\ldots,Y_{10}]
\]
be the Macaulay inverse system.  It is a contraction-stable
\(S_0\)-module of dimension \(584\).

\begin{theorem}[Minimal inverse-system generators]
\label{thm:inverse}
The quotient \(D/\mathfrak mD\) has dimension \(60\).  Its top-degree and
torus-refined generating polynomial is
\begin{align*}
\mathcal S(t,z)={}&t^5(z^{-1}+z^3)\\
&+t^6(z^{-4}+z^{-3}+5z^{-2}+11z^{-1}+8+6z+z^2)\\
&+t^7(z^{-4}+2z^{-3}+3z^{-2}+7z^{-1}+3+3z+2z^2+z^3)\\
&+t^8(z^{-1}+1+z).
\end{align*}
\end{theorem}

\begin{proof}[Computer-assisted proof]
Modulo \(p=1{,}000{,}003\), the weight-block dimensions of \(D\), of
\(\mathfrak mD\), and of their quotient are as follows:
\[
\begin{gathered}
\begin{array}{c|rrrrrrr}
\text{weight}&-7&-6&-5&-4&-3&-2&-1\\ \hline
\dim D_w&1&4&12&28&47&74&94\\
\dim(\mathfrak mD)_w&1&4&12&26&44&66&74\\
\dim\text{ quotient}&0&0&0&2&3&8&20
\end{array}\\[6pt]
\begin{array}{c|rrrrrr}
\text{weight}&0&1&2&3&4&5\\ \hline
\dim D_w&96&86&67&45&23&7\\
\dim(\mathfrak mD)_w&84&76&64&43&23&7\\
\dim\text{ quotient}&12&10&3&2&0&0.
\end{array}
\end{gathered}
\]
The quotient dimensions sum to 60, so specialization gives
\(\dim_\mathbb Q D/\mathfrak mD\le60\).

The certificate contains 60 exact rational divided-power polynomials with
the filtered character displayed in the theorem.  Their reductions form a
basis of the modular quotient, so their characteristic-zero classes are
independent and give the reverse inequality.  Their complete contraction
closure modulo \(p\) has rank \(584\): 60 independent starting vectors
produce 465 new independent first contractions and 59 new independent second
contractions.  Nakayama's lemma therefore proves that the 60 rational
elements generate \(D\).
\end{proof}

\begin{corollary}[Socle]
\label{cor:socle}
The Cohen--Macaulay type is
\[
\dim_\mathbb Q\operatorname{Soc}(R)=60.
\]
The associated-graded socle dimensions occur as
\[
\dim\frac{\operatorname{Soc}(R)\cap\mathfrak m^d}
{\operatorname{Soc}(R)\cap\mathfrak m^{d+1}}
=
\begin{cases}
2,&d=5,\\
33,&d=6,\\
22,&d=7,\\
3,&d=8,\\
0,&\text{otherwise}.
\end{cases}
\]
In particular, \(R\) is neither Gorenstein nor level.
\end{corollary}

\begin{proof}
Macaulay duality gives a perfect pairing
\[
\operatorname{Soc}(R)\times D/\mathfrak mD\longrightarrow\mathbb Q.
\]
The filtered statement follows from
\[
F_dD=(\mathfrak m^{d+1})^\perp
\]
and the top-degree distribution in \cref{thm:inverse}.  Type 60 excludes
Gorensteinness, while the occurrence of socle classes in four distinct
\(\mathfrak m\)-adic degrees excludes levelness.
\end{proof}

\section{Minimal Kuranishi equations}
\label{sec:equations}

\begin{theorem}[Thirty-six minimal equations]
\label{thm:equations}
The conormal space of the Kuranishi ideal has dimension
\[
\mu(I_\kappa)
=\dim_\mathbb Q I_\kappa/\mathfrak mI_\kappa=36.
\]
Its initial-order distribution is
\[
11,\ 13,\ 11,\ 0,\ 1
\]
in orders \(2,3,4,5,6\), respectively.  The corresponding torus characters
are
\begin{align*}
E_2(z)={}&z^{-6}+z^{-5}+z^{-2}+z+z^2+z^3+2z^4+2z^5+z^6,\\
E_3(z)={}&z^{-3}+2z^{-2}+2z^{-1}+1+2z+z^2+z^3+z^5+2z^6,\\
E_4(z)={}&z^{-8}+z^{-7}+2z^{-6}+2z^{-5}+z^{-4}+z^{-3}+z^{-2}+2z^5,\\
E_6(z)={}&z^3.
\end{align*}
No new minimal generator first appears in orders seven, eight, or nine.
\end{theorem}

\begin{proof}[Computer-assisted proof]
Exact filtered Kuranishi elimination produces 36 independent classes with the
displayed orders and weights.  Independently within the computational
pipeline, the Koszul complex of the modular algebra gives
\[
\dim_{\mathbb F_p}H_1(u_1,\ldots,u_{10};R_{\mathbb F_p})=36
\]
with the same weight character.  Since
\[
H_1(u;R)\cong
\operatorname{Tor}_1^{S_0}(R,\mathbb Q)
\cong I_\kappa/\mathfrak mI_\kappa,
\]
specialization gives the matching characteristic-zero upper bound.  The 36
exact classes give the lower bound.
\end{proof}

\begin{corollary}[Extremal Betti numbers]
\label{cor:betti}
For the minimal \(S_0\)-free resolution of \(R\),
\[
\beta_1^{S_0}(R)=36,\qquad
\beta_{10}^{S_0}(R)=60.
\]
The ring is not a complete intersection.
\end{corollary}

\begin{proof}
The first equality is \cref{thm:equations}.  Since \(R\) is Artinian and
\(S_0\) is regular local of dimension ten, local duality identifies the last
Betti number with the type, which is 60 by \cref{cor:socle}.  A
zero-dimensional complete intersection in a ten-dimensional regular local
ring has ten minimal defining equations, not 36.
\end{proof}

\section{Exact computation and reproducibility}
\label{sec:computation}

All defining coefficients are rational.  No floating-point arithmetic is
used.  The proof is divided into small exact certificates rather than one
uninspectable final assertion.

The certificate layers are:
\begin{enumerate}[label=(\arabic*)]
\item stabilizer and residual-parameter identities, including the torus
\eqref{eq:torus} and the omitted curve;
\item the slice rank, ten tangent vectors and weights, and the quadratic
obstruction space;
\item the positive, negative, and fixed-locus calculations of
\cref{prop:pure-weight};
\item 7,300 targeted rational Kuranishi coefficients through order eight and
comparison of all 990 new coefficients with the full modular output;
\item fresh CRT reconstruction and exact all-row checks for the three
terminal dual classes;
\item the nine order-nine determinants in \cref{prop:order-nine};
\item the contraction quotient and closure, giving 60 minimal
inverse-system generators and rank \(584\); and
\item the Koszul \(H_1\) calculation giving 36 minimal equations with their
weight character.
\end{enumerate}

The fresh audit reran the original SHA-256 manifests before executing the
programs.  Generated reduced-rigidity files agreed byte-for-byte with the
recovered outputs.  The order-eight classes were reconstructed afresh; the
three runs required 17, 24, and 31 small primes and reproduced the archived
rational supports and coefficients.  The order-nine JSON and text
certificates also agreed byte-for-byte.  Finally, the three self-contained
inverse-system verifiers reproduced type 60, closure rank 584, and Koszul
\(H_1\) dimension 36.

\subsection*{What this verification does not establish}

The several checks are algebraically distinct, but they are not a
genuinely independent implementation: the principal calculation and its
verifiers belong to one Python/SymPy-centered family and share some exported
data.  They verify exact identities, ranks, specializations, and
reconstructions.  They do not independently validate the geometric choice of
moduli problem, the interpretation of the slice, or the novelty claim.

The computational supplement now contains the exact Kuranishi rows, the 60
rational inverse-system generators, the order-eight and order-nine
certificates, the border-basis reconstruction data, every verifier,
dependency versions, source hashes, and fresh logs.  A compact index maps
claims to source files.  The large exact matrices and modular-reconstruction
files are shipped in the companion data archive rather than embedded in the
paper source.  This makes the number \(584\) inspectable in two complementary
ways: as the closure of the inverse system under contraction and as the
dimension of the commuting multiplication representation.  A second
implementation in Singular or Macaulay2 remains desirable independent
confirmation, not a missing component of the stated exact certificate.

\section{Interpretation and open directions}
\label{sec:interpretation}

The reduced answer is simple and the scheme-theoretic answer is not.  The
degree-seven counterexample is isolated in the stated reduced affine-moduli
germ, yet it sits on a length-\(584\) infinitesimal thickening.  The ring has
embedding dimension ten, 36 minimal equations, type 60, and socle in four
different orders.  Thus it is not well modeled by a hypersurface, complete
intersection, Gorenstein algebra, or even a level algebra.

The number \(584\) does not count nearby maps.  It is the intersection
multiplicity created when the formally smooth source-reparametrization orbit
meets the degree-seven coefficient stage.  A one-variable model is
\[
V(y-x^9)\cap V(y)
=\operatorname{Spec}\C[[x]]/(x^9):
\]
the reduced intersection is one point, although the two conditions have
ninth-order contact.  The algebra \(\mathcal R\) is a ten-variable,
non-complete-intersection analogue.  Its Hilbert function records the
successive layers of the transverse normal cone, not a collection of
distinct counterexamples.

Three questions are deliberately left outside the theorem.

\begin{question}
Can the complete local algebra be recovered conceptually from the
marked-root geometry, rather than by coefficient elimination?
\end{question}

\begin{question}
What is the reduced deformation germ when ordinary degree eight is allowed?
Existing calculations give a mixture of exact and modular finite-order
evidence, but not a characteristic-zero radical theorem.
\end{question}

\begin{question}
How does the Artin algebra change under polynomial, rather than affine,
left--right equivalence or after stabilization?
\end{question}

\begin{question}[Independent source-flow verification]
\label{q:source-flow-verifier}
Can one derive the finite overflow and Kuranishi equations directly from
volume-preserving flows on the source?  In view of
\cref{prop:root-coordinate-divergence}, the sharper target is an independent
construction using the different filtration or transferred
\(L_\infty\)-data.  Agreement would test the source-orbit interpretation of
the local Artin algebra.
\end{question}

\appendix
\section{Root coordinates and the source-flow complex}
\label{app:root-coordinate-source-flow}

The marked-root coordinates make the volume-preserving source action much
sparser than it appears in coefficient coordinates.  Write the target cubic
as
\[
As^3+Bs^2+bs-2a=0
\]
and use \((A,B,s)\) as source coordinates.  Solving for the remaining source
coordinates in the standard marked-root frame gives the Jacobian factor
\[
\delta=1-Bs+3As^2.
\]

\begin{proposition}[Weighted divergence in root coordinates]
\label{prop:root-coordinate-divergence}
In the marked-root chart,
\[
dx\wedge dy\wedge dz
=-\delta\,dA\wedge dB\wedge ds.
\]
Consequently a vector field
\[
X=a\,\partial_A+b\,\partial_B+c\,\partial_s
\]
preserves source volume if and only if
\[
\partial_A(\delta a)+\partial_B(\delta b)+\partial_s(\delta c)=0.
\]
\end{proposition}

\begin{proof}
The first identity is the determinant of the explicit triangular change
from \((x,y,z)\) to the marked-root variables.  Cartan's formula gives
\[
\mathcal L_X(\delta\,dA\wedge dB\wedge ds)
=
\bigl(
\partial_A(\delta a)+\partial_B(\delta b)+\partial_s(\delta c)
\bigr)dA\wedge dB\wedge ds,
\]
which proves the second assertion.
\end{proof}

The same factor \(\delta\) is the derivative of the cubic in the marked
root.  Thus the different of the cubic cover is exactly the weight in the
source divergence equation.  This suggests filtering the volume-preserving
Lie algebra by powers of the different, transferring it to the finite
overflow complex, and comparing the resulting higher brackets with the
Kuranishi equations computed in the body.

\begin{question}[Different-filtered reconstruction]
\label{q:different-filtered-reconstruction}
Can a different-filtered source-flow complex, or an equivalent transferred
\(L_\infty\)-model, recover the bounded-degree Kuranishi ideal and its
length \(584\) independently of coefficient elimination?
\end{question}

\section{Independent source-flow reconstruction through order five}
\label{app:source-flow-reconstruction-progress}

Marked-root coordinates give a second construction of the bounded-degree
deformation equations.  This appendix records the part reconstructed
independently of the archived coefficient-space matrices.

\subsection{Root coordinates and the exact degree filtration}

For the unnormalized target \(F=(a,b,c)=(P,Q,R)\), put
\[
r=\frac2x,\qquad t=y+\frac1x,\qquad c=R.
\]
Then
\[
a=t^2+\frac12rt-ct^3,\qquad
b=r+4t-3ct^2,
\]
and
\[
x=\frac2r,\qquad y=t-\frac r2,\qquad
z=\frac{10r^2-12tr-cr^3}{8}.
\]
The source volume form is
\[
dx\wedge dy\wedge dz=-\frac r4\,dc\wedge dt\wedge dr.
\]
Thus \(X=A\partial_c+B\partial_t+C\partial_r\) is volume preserving exactly
when
\[
\partial_c(rA)+\partial_t(rB)+\partial_r(rC)=0.
\]
Ordinary degree is retained without approximation by substituting the
ordinary monomial basis, or equivalently by using the Rees generators
\[
qx=\frac{2q}{r},\qquad
qy=q\left(t-\frac r2\right),\qquad
qz=\frac q8(10r^2-12tr-cr^3).
\]

\begin{theorem}[Source-flow reconstruction]
\label{thm:source-flow-reconstruction}
The direct determinant complex and the marked-root weighted-divergence
complex give the same filtered transverse Kuranishi ideal through parameter
order four over \(\mathbb Q\).  In particular they give
\[
H(0),\ldots,H(4)=(1,10,44,108,157)
\]
and filtered minimal-generator counts
\[
11\text{ quadratic},\qquad
13\text{ cubic},\qquad
11\text{ quartic}.
\]
\end{theorem}

\begin{proof}[Exact computer-assisted proof]
The two implementations build their linear maps independently.  The direct
model differentiates the determinant in the 337-dimensional affine slice.
The root-native model contracts the weighted-divergence complex inside the
exact degree-seven Rees basis.  Exact rational row reduction gives rank
\(327\) and the same ten tangent weights
\[
(-1,2,-3,-2,-1,0,1,1,2,3).
\]
At each later order a contracting homotopy is chosen inside the same model.
The canonical rational equations and filtered Macaulay row spaces then agree
coefficientwise through order four.
\end{proof}

The root-native contraction is essential.  Importing the direct model's
chosen quadratic and cubic corrections into the divergence complex gives a
different quartic section and an incorrect rank.  Once both homotopies are
chosen natively, their obstruction spaces agree.

There is also a genuine filtered-versus-graded distinction.  The associated
graded ideal has twelve new quartic initial classes, but the filtered ideal
has only eleven minimal generators of initial order four.  The extra initial
class is created by cancellation among lower-order equations and their
higher tails.

\subsection{Independent reduced rigidity}

\begin{corollary}[Source-flow proof of reduced affine rigidity]
\label{cor:source-flow-reduced-rigidity}
In the specified degree-seven affine slice,
\[
\sqrt I=(u_1,\ldots,u_{10}).
\]
\end{corollary}

\begin{proof}
The reconstructed positive-weight equations have unit ideal on the only
surviving weight-one chart, with the final two residual quadratics having
nonzero exact resultant.  On the negative stratum, equations through order
four give
\[
\sqrt{I^-}=(u_1,u_3,u_4,u_5);
\]
the needed nilpotent powers are certified by exact ideal-membership
identities.  On the fixed stratum, the remaining equation is \(u_6^3\)
times a formal unit.  The torus-nullcone lemma then gives the stated radical.
\end{proof}

\subsection{Order five and the sextic frontier}

\begin{proposition}[Exact fifth-order calculation]
\label{prop:source-flow-order-five}
At parameter order five the direct model has Macaulay rank \(1857\), so
\[
H(5)=2002-1857=145.
\]
There is no new minimal quintic generator.
\end{proposition}

\begin{proof}[Exact computer-assisted proof]
Every torus-weight block has a rational rank certificate.  Modular
elimination selects a full-rank minor, while exact rational solves express
all remaining columns in its span.  A second calculation of
\(I/\mathfrak m I\) shows that all fifth-order equations are tails of the
existing \(35\) filtered generators.
\end{proof}

At order six, two independent good primes give the same blockwise result
\[
H(6)=86
\]
and one new minimal class of torus weight \(3\), in agreement with the
sextic generator in the complete coefficient-space computation.  The
recovered exact run log covers the weight blocks from \(-18\) through \(0\).
For each weight from \(6\) through \(18\), a newly generated nonzero maximal
minor modulo a good prime proves that the rational Macaulay block has full
column rank; the calculation was repeated at two primes.  For the deficient
weights \(1,\ldots,5\), exact rational row reduction gives column relations
that were checked against every original row, completing the rank proof in
characteristic zero.

The same method applied to \(I/\mathfrak mI\) at weight \(3\) gives
\[
\operatorname{rank}(\mathfrak mI)=542,\qquad
\operatorname{rank}(I)=545,
\]
while their pure sextic initial spaces have dimensions \(341\) and \(342\).
Consequently there is exactly one new sextic class.  With variables ordered
as in the transverse slice, one primitive integer representative is
\[
642816u_1u_6u_7^4-60u_4u_7u_8^4+5u_4u_8^5
-75u_5u_6u_8^4.
\]
This completes the independently reconstructed direct-model calculation
through order six over \(\mathbb Q\).  It does not extend the exact
root/Rees row-space comparison beyond order four, nor does it reconstruct
orders seven and eight.

The reconstructed source-flow results establish the geometric meaning of
the linear-through-quintic deformation data and independently prove the
reduced germ.  They do not yet reconstruct the complete length-\(584\)
nonreduced algebra from the source-flow complex.

\input{appendices/degree-eight}
\input{appendices/border-basis-and-betti}
\section*{Companion results and research register}

This reader edition is accompanied by a separate \emph{Results and Research
Register}.  The register preserves secondary proved results, conditional
developments, open questions, corrections, and computational leads without
making them part of this paper's theorem spine.  Proof-bearing technical
developments omitted from this reader PDF remain in the versioned source
tree and in the complete version-8 archival edition.  The v9 primary-home
map distinguishes the paper responsible for a result from other papers that
use or motivate it.


\section*{AI and computational disclosure}

AI systems were used extensively in exploration, symbolic-program
development, proof drafting, source organization, and manuscript editing.
They are not authors.  The mathematical evidence for the computer-assisted
theorems consists of the exact certificates and reproducible calculations
described above, not the output of a language model.  Human specialist review
has not yet been completed.

\bibliographystyle{amsplain}
\bibliography{../common/references}

\end{document}
